\label{sec:vm-deployment}

Für den Produktivbetrieb läuft LearnHub auf einer virtuellen Maschine 
der HTL Leonding unter der Adresse \texttt{vm91.htl-leonding.ac.at}. 
Die VM läuft mit Ubuntu Server und hat Docker Engine sowie Docker 
Compose installiert. Der Vorteil gegenüber einem direkten Deployment 
auf physischer Hardware ist, dass die VM unabhängig von lokalen 
Entwicklerrechnern betrieben, gesichert und administriert werden 
kann \cite{pahl2015containerization}. Da alle Services in Containern 
laufen, muss auf der VM selbst keine zusätzliche Software installiert 
werden, also kein Java, kein Node.js und keine Datenbank. Alles was 
die Anwendung braucht ist in den Images enthalten.

\subsubsection{Einrichtung der VM}

Nach der Installation von Docker auf der VM wird als erstes das 
externe Docker-Netzwerk \texttt{learnhubnet} vom Traefik Stack 
angelegt. Dieses Netzwerk muss existieren bevor der LearnHub Stack 
gestartet werden kann, da beide Stacks darüber kommunizieren. 
Danach werden die beiden Compose Files auf die VM übertragen. 
Zuerst wird der Traefik Stack gestartet, der das Netzwerk erstellt 
und den Reverse-Proxy hochfährt. Erst danach wird der LearnHub 
Stack gestartet, dessen Services sich automatisch in das bestehende 
Netzwerk einklinken und ihre Routing-Regeln an Traefik melden.

Für HTTPS wird ein SSL-Zertifikat benötigt. Traefik übernimmt 
dabei die automatische Ausstellung und Erneuerung des Zertifikats 
über Let's Encrypt. Die Zertifikatsdaten werden im \texttt{certs} 
Ordner auf der VM gespeichert, der als Volume in den Traefik 
Container gemountet ist. Dadurch bleibt das Zertifikat auch nach 
einem Neustart des Containers erhalten.

\subsubsection{Deployment-Prozess}

Der Deployment-Prozess läuft in wenigen klar definierten Schritten ab:

\begin{enumerate}
    \item Die Docker-Images für Frontend und Backend werden lokal 
    gebaut und auf Docker Hub hochgeladen, wo sie öffentlich 
    verfügbar sind.
    \item Die Compose Files werden auf die VM übertragen.
    \item Der Traefik Stack wird gestartet: 
    \begin{verbatim}
docker compose up -d
    \end{verbatim}
    \item Der LearnHub Stack wird gestartet:
    \begin{verbatim}
docker compose up -d
    \end{verbatim}
    \item MySQL führt seinen Healthcheck durch. Sobald die Datenbank 
    bereit ist, startet das Backend automatisch und baut die 
    Datenbankverbindung auf.
\end{enumerate}

Nach dem Start ist die Anwendung unter 
\texttt{https://vm91.htl-leonding.ac.at} erreichbar. Alle Container 
sind mit \texttt{restart: always} konfiguriert, was bedeutet dass 
sie bei einem Absturz oder Neustart der VM automatisch wieder 
hochfahren, ohne dass manuell eingegriffen werden muss.

\subsubsection{Updates}

Um eine neue Version der Anwendung zu deployen, wird das neue 
Docker-Image gebaut und auf Docker Hub hochgeladen. Auf der VM 
reichen dann zwei Befehle:

\begin{verbatim}
docker compose pull
docker compose up -d
\end{verbatim}

\texttt{docker compose pull} lädt die neuesten Images von Docker 
Hub herunter. \texttt{docker compose up -d} startet danach nur 
jene Container neu, deren Image sich geändert hat. Services die 
nicht aktualisiert wurden laufen ohne Unterbrechung weiter. Da 
die Daten in Volumes gespeichert sind, gehen bei einem Update 
weder Datenbankeinträge noch hochgeladene Dateien verloren. Das 
macht Updates unkompliziert und sicher, da im schlimmsten Fall 
einfach auf das vorherige Image zurückgewechselt werden kann.

\subsubsection{Monitoring}

Traefik bietet ein eingebautes Dashboard das unter Port 8080 
erreichbar ist. Dort ist auf einen Blick sichtbar welche Services 
gerade laufen, welche Routen aktiv sind und ob alle Container 
erreichbar sind. Für eine detailliertere Übersicht können auf 
der VM direkt Docker-Befehle verwendet werden:

\begin{verbatim}
docker compose ps
docker compose logs -f
\end{verbatim}

\texttt{docker compose ps} zeigt den Status aller Container, 
\texttt{docker compose logs -f} streamt die Logs aller Services 
in Echtzeit. Das ist besonders hilfreich um Fehler beim Start 
der Anwendung oder im laufenden Betrieb schnell zu finden.